\chapter{Measurement, Unitary Evolution, and Quantum Dynamics}

\lectureoverview{We first collect the measurement principles already in use: observables are Hermitian operators, their eigenvalues are possible outcomes, distinguishable states are orthogonal, and Born's rule assigns probabilities. The new problem is dynamics. Reversibility requires a common linear evolution that preserves overlaps; this makes the evolution operator unitary and its infinitesimal generator Hermitian. The generator is the Hamiltonian, and its action gives the time-dependent Schr\"odinger equation. We then derive the expectation-value formula, the commutator equation for averages, its relation to classical Poisson brackets, and the conservation of average energy.}

For the moment, put the spin system aside. We will return to it after the general structure is in place. The principles below are not presented as a minimal independent set. Some reinforce or partly imply others, but stating them separately makes their physical content easier to see.

\section{What a Measurement Can Return}

An \emph{observable} is something that can be measured. The established word is slightly grander than necessary; \emph{measurable} would have served just as well. To every observable we associate a Hermitian operator. Let \(L\) denote such an operator for now:
\begin{equation}
L^\dagger=L.
\end{equation}
Not every linear operator is Hermitian. Conversely, whether a particular Hermitian operator can be measured with available materials and apparatus may be a difficult experimental question. The abstract working rule is nevertheless generous: a Hermitian operator represents a possible observable, even if the practical measurement has not yet been invented.

If
\begin{equation}
L\ket{\lambda}=\lambda\ket{\lambda},
\end{equation}
then \(\lambda\) is a possible numerical result of measuring \(L\). We use the same symbol to label the eigenvalue and its eigenvector. The ket bars distinguish the two: \(\lambda\) is a number, while \(\ket{\lambda}\) is a state.

Why Hermitian operators? A primitive apparatus reports real numbers. If an instrument appears to report one complex number, it is really reporting two independent real quantities. Thus observable eigenvalues must be real. Moreover, states associated with different definite outcomes must be physically distinguishable and hence orthogonal. In the finite-dimensional setting used here, a linear operator with a complete orthonormal eigenbasis and real eigenvalues is Hermitian. The physical statements and the Hermitian structure therefore fit one another.

\subsection{Distinguishability Is Orthogonality}

The phrase \emph{physically distinguishable} has an operational meaning. Suppose a spin is prepared either up or down. A measurement along the \(z\)-axis separates the possibilities without ambiguity: up gives \(+1\), down gives \(-1\). Likewise, left and right can be distinguished by measuring along the \(x\)-axis.

Now suppose the preparation is either up or right. A \(z\)-measurement gives \(+1\) with certainty for up, but the right state also gives \(+1\) with nonzero probability. An \(x\)-measurement merely reverses the asymmetry. Tilting the apparatus by \(45^\circ\) does not remove the overlap between the outcome distributions. No orientation produces a result that always distinguishes the two preparations.

\begin{definition}
Two states are physically distinguishable when some experiment, or collection of experiments, can determine unambiguously which preparation was supplied.
\end{definition}

Quantum mechanics represents complete distinguishability geometrically:
\begin{equation}
\text{distinguishable states}
\quad\Longrightarrow\quad
\text{orthogonal vectors}.
\end{equation}

\begin{classroomqa}
\audiencequestion{Could a different measurement axis distinguish an up state from a right state?}
\lecturerresponse{No. Measuring along \(z\) or \(x\) can sometimes identify the preparation, but at least one result remains compatible with both. An intermediate axis also leaves overlapping probabilities. The obstruction is not a poor choice of apparatus; the two state vectors are not orthogonal.}
\end{classroomqa}

\section{Born's Rule}

The eigenvalues tell us what may happen, but not how often. The probability rule supplies the predictive content. If a normalized system is prepared in \(\ket{A}\), and a normalized eigenvector \(\ket{\lambda}\) belongs to the measured observable \(L\), then
\begin{equation}
P(\lambda)=\left|\braket{\lambda}{A}\right|^2.
\label{eq:l4-born-rule}
\end{equation}
The overlap \(\braket{\lambda}{A}\) is a \emph{probability amplitude}. It is generally complex, so it cannot itself be a probability. Its absolute square is real and nonnegative.

\begin{figure}[htbp]
\centering
\includegraphics[width=0.86\textwidth]{lecture_04_figure_01.png}
\caption{Born's rule written beside the measurement principles: probability is the absolute square of an overlap.}
\lectureframe{00:12:40}
\end{figure}

Choose an orthonormal eigenbasis of \(L\) and expand
\begin{equation}
\ket{A}=\sum_i \alpha_i\ket{\lambda_i},
\qquad
\alpha_i=\braket{\lambda_i}{A}.
\end{equation}
Born's rule becomes
\begin{equation}
P(\lambda_i)=|\alpha_i|^2.
\end{equation}
Normalization then gives
\begin{equation}
\sum_i P(\lambda_i)=\sum_i|\alpha_i|^2=1.
\end{equation}
The positivity and normalization requirements do not by themselves derive Born's rule, but they show why the absolute square fits probability theory so naturally. Ultimately the rule is tested by experiment.

\begin{classroomqa}
\audiencequestion{Is the absolute square also required so that all probabilities add to one?}
\lecturerresponse{For a normalized state expanded in an orthonormal basis, the squared amplitudes do add to one. Positivity and normalization therefore fit Born's rule perfectly. Which statement should be regarded as logically prior is less important here than their mutual consistency; experiment fixes the probability rule.}
\end{classroomqa}

\begin{classroomqa}
\audiencequestion{Why is an eigenvalue, which is a number, written inside a ket?}
\lecturerresponse{The notation \(\ket{\lambda}\) means the eigenvector associated with the numerical eigenvalue \(\lambda\). A more punctilious notation could introduce another label, but too many indices soon obscure the calculation. The ket symbol keeps the number and vector distinct.}
\end{classroomqa}

\begin{classroomqa}
\audiencequestion{Which principle turns several possible outcomes into one actual outcome?}
\lecturerresponse{Born's rule assigns the probabilities. If the prepared state is already an eigenvector of \(L\), its corresponding outcome has probability one. For a superposition, one outcome is obtained with the probabilities in Equation~\eqref{eq:l4-born-rule}. The associated state-update or collapse rule is a further measurement question; it will be revisited after entanglement is available.}
\end{classroomqa}

\subsection{Degenerate Outcomes}

An eigenvalue may belong to more than one independent eigenvector. Suppose
\begin{align}
L\ket{\lambda,1}&=\lambda\ket{\lambda,1},\\
L\ket{\lambda,2}&=\lambda\ket{\lambda,2}.
\end{align}
Every linear combination in their span is also an eigenvector with eigenvalue \(\lambda\):
\begin{equation}
L\bigl(\alpha\ket{\lambda,1}+\beta\ket{\lambda,2}\bigr)
=
\lambda\bigl(\alpha\ket{\lambda,1}+\beta\ket{\lambda,2}\bigr).
\end{equation}
This freedom lets us choose an orthonormal basis inside the degenerate eigenspace. If the label \(a\) runs over that basis, the probability of the numerical outcome \(\lambda\) is the sum over its mutually exclusive orthogonal channels:
\begin{equation}
P(\lambda)=\sum_a\left|\braket{\lambda,a}{A}\right|^2.
\label{eq:l4-degenerate-born}
\end{equation}
Equivalently, if \(P_\lambda\) is the projector onto the full eigenspace, then \(P(\lambda)=\bra{A}P_\lambda\ket{A}\). The projector notation is compact; Equation~\eqref{eq:l4-degenerate-born} displays the probability addition directly.

\begin{classroomqa}
\audiencequestion{What if the prepared state overlaps several eigenvectors, or several eigenvectors share the same eigenvalue?}
\lecturerresponse{For different eigenvalues, each overlap supplies the probability of its own outcome. For several orthogonal eigenvectors with the same eigenvalue, square the overlaps and add them. If the vectors originally chosen in the degenerate subspace are not orthogonal, replace them by an orthonormal basis first; the summed probability is basis independent.}
\end{classroomqa}

\begin{classroomqa}
\audiencequestion{Can one obtain an unambiguous answer from a spin whose preparation is completely unknown?}
\lecturerresponse{Not from one arbitrary specimen. Without restrictions on how it was prepared, there is no single measurement that identifies an unknown quantum state unambiguously. Preparation information specifies the alternatives among which a measurement can discriminate.}
\end{classroomqa}

\section{The New Question: How States Change}

The preceding principles specify states and measurements. Dynamics asks what happens between one time and another. The classical analogy is useful. A two-state coin might remain unchanged at every tick, or it might alternate heads, tails, heads, tails. A six-sided die admits more update rules, and a continuous state space admits infinitely many. A dynamical law maps the present state to the next one.

A satisfactory classical law is deterministic and reversible. Distinct states do not merge into one state, because if they did the present would no longer determine the past. This preservation of information is sometimes called the \emph{minus first law}: it comes before the more familiar conservation laws.

The quantum version begins with a state vector \(\ket{\psi(t)}\), explicitly regarded as a function of time. Postulate a single linear time-development operator \(U(t)\) such that
\begin{equation}
\ket{\psi(t)}=U(t)\ket{\psi(0)}.
\label{eq:l4-time-development}
\end{equation}
The same \(U(t)\) acts on every initial state of the system. It is not chosen afresh according to the input vector, and its action is linear. More general nonlinear or state-dependent laws can be imagined, but they rapidly conflict with the probabilistic, causal, and composite-system structure of quantum mechanics.

\subsection{Reversibility Preserves Overlaps}

Let a second initial state evolve by the same law:
\begin{equation}
\ket{\phi(t)}=U(t)\ket{\phi(0)}.
\end{equation}
If the two preparations are distinguishable initially, they must remain distinguishable. Therefore
\begin{equation}
\braket{\psi(0)}{\phi(0)}=0
\quad\Longrightarrow\quad
\braket{\psi(t)}{\phi(t)}=0.
\end{equation}

\begin{figure}[htbp]
\centering
\includegraphics[width=0.86\textwidth]{lecture_04_figure_02.png}
\caption{Two states evolve under the same \(U(t)\); initial orthogonality remains orthogonality.}
\lectureframe{00:34:43}
\end{figure}

Preserving normalized states and orthogonality is strong enough to preserve all inner products. To see the mechanism, expand two arbitrary vectors in an orthonormal basis. The basis vectors remain orthonormal under the evolution, so the coefficients reconstruct the same overlap at every time:
\begin{equation}
\braket{\psi(t)}{\phi(t)}=\braket{\psi(0)}{\phi(0)}.
\label{eq:l4-overlap-conservation}
\end{equation}
Thus not only complete distinguishability but the entire degree of similarity between two states is conserved.

The bra corresponding to Equation~\eqref{eq:l4-time-development} is obtained by Hermitian conjugation:
\begin{equation}
\bra{\psi(t)}=\bra{\psi(0)}U^\dagger(t).
\end{equation}
Substituting the bra and ket evolutions into Equation~\eqref{eq:l4-overlap-conservation} gives
\begin{equation}
\bra{\psi(0)}U^\dagger(t)U(t)\ket{\phi(0)}
=
\braket{\psi(0)}{\phi(0)}
\end{equation}
for every pair of initial vectors. If an operator has the same matrix elements as the identity between every pair of vectors, it is the identity. Hence
\begin{equation}
U^\dagger(t)U(t)=I.
\label{eq:l4-unitarity}
\end{equation}
An operator satisfying Equation~\eqref{eq:l4-unitarity} is \emph{unitary}. In the finite-dimensional setting,
\begin{equation}
U^\dagger(t)=U^{-1}(t),
\end{equation}
so the evolution can be run backward.

\begin{classroomqa}
\audiencequestion{Do uppercase and lowercase \(\Phi\) denote different kinds of states here?}
\lecturerresponse{No. That is merely inconsistent handwriting. We use \(\ket{\phi}\) consistently; the mathematics distinguishes states by their labels and values, not by accidental capitalization.}
\end{classroomqa}

\begin{classroomqa}
\audiencequestion{Can \(U^\dagger U=I\) be read as saying that one factor is the inverse of the other?}
\lecturerresponse{Yes. A unitary operator has inverse \(U^{-1}=U^\dagger\). This is the algebraic expression of reversible, overlap-preserving evolution.}
\end{classroomqa}

\section{Infinitesimal Time and the Hamiltonian}

At zero elapsed time nothing changes, so \(U(0)=I\). For a very small interval \(\epsilon\), \(U(\epsilon)\) must be close to the identity. In units that temporarily set \(\hbar=1\), write
\begin{equation}
U(\epsilon)=I-i\epsilon H+O(\epsilon^2).
\label{eq:l4-small-u}
\end{equation}
The minus sign and factor of \(i\) are initially conventions; without knowing \(H\), they can be absorbed into its definition. Their usefulness appears when unitarity is imposed. Hermitian conjugation gives
\begin{equation}
U^\dagger(\epsilon)=I+i\epsilon H^\dagger+O(\epsilon^2).
\end{equation}
Then
\begin{align}
I
&=U^\dagger(\epsilon)U(\epsilon)\\
&=\bigl(I+i\epsilon H^\dagger\bigr)
  \bigl(I-i\epsilon H\bigr)+O(\epsilon^2)\\
&=I+i\epsilon(H^\dagger-H)+O(\epsilon^2).
\end{align}
The first-order term must vanish:
\begin{equation}
H^\dagger=H.
\label{eq:l4-h-hermitian}
\end{equation}

\begin{figure}[htbp]
\centering
\includegraphics[width=0.86\textwidth]{lecture_04_figure_03.png}
\caption{The first-order unitarity calculation forces the infinitesimal generator to be Hermitian.}
\lectureframe{00:43:38}
\end{figure}

The letter \(H\) stands not for Hermitian but for \emph{Hamiltonian}. Since it is Hermitian, it represents an observable. It is identified with the energy and, more structurally, with the generator of time translations. The identification is not exhausted by a name: in systems with a known classical limit, the quantum \(H\) must reproduce the classical Hamiltonian's dynamics. Symmetry, the degrees of freedom, physical judgment, and ultimately experiment determine which Hermitian operator is the Hamiltonian of a particular system.

\begin{classroomqa}
\audiencequestion{Could any Hermitian operator be used for \(H\)? What selects the actual one?}
\lecturerresponse{The derivation so far allows any Hermitian operator on the system's Hilbert space. Physics selects one built from that system's degrees of freedom. For a spin, it must be constructed from spin operators and external fields. Symmetry and known interactions narrow the possibilities; experiment fixes the remaining constants and sometimes the form itself.}
\end{classroomqa}

The Hamiltonian may depend explicitly on time. A magnetic field whose current is being varied supplies a simple example. If the external parameters are time independent, then \(H\) has no explicit time dependence; this is time-translation invariance. If the parameters are driven in time, the energy need not be conserved.

\begin{classroomqa}
\audiencequestion{In the small-time expansion, is \(H\) a constant, or should it be \(H(0)\)?}
\lecturerresponse{For a time-independent problem, \(H\) is the same operator at every time. More generally, the local step beginning at time \(t\) contains \(H(t)\). The derivation here first treats the time-translation-invariant case and later notes how explicit time dependence changes energy conservation.}
\end{classroomqa}

\section{The Time-Dependent Schr\"odinger Equation}

Apply Equation~\eqref{eq:l4-small-u} to the state:
\begin{equation}
\ket{\psi(\epsilon)}=(I-i\epsilon H)\ket{\psi(0)}+O(\epsilon^2).
\end{equation}
Rearranging and dividing by \(\epsilon\),
\begin{equation}
\frac{\ket{\psi(\epsilon)}-\ket{\psi(0)}}{\epsilon}
=-iH\ket{\psi(0)}+O(\epsilon).
\label{eq:l4-difference-quotient}
\end{equation}

\begin{figure}[htbp]
\centering
\includegraphics[width=0.86\textwidth]{lecture_04_figure_04.png}
\caption{The infinitesimal state update rearranged into a vector-valued difference quotient.}
\lectureframe{00:51:42}
\end{figure}

There is no special mystery in differentiating a vector. Two vectors can be subtracted, their difference can be divided by a number, and the limit defines another vector. The reference time called zero above can be any instant. Taking \(\epsilon\to0\) therefore yields
\begin{equation}
\boxed{\frac{d}{dt}\ket{\psi(t)}=-iH\ket{\psi(t)}}
\qquad(\hbar=1).
\label{eq:l4-schrodinger-natural}
\end{equation}

\begin{figure}[htbp]
\centering
\includegraphics[width=0.86\textwidth]{lecture_04_figure_05.png}
\caption{The completed board derivation culminates in the boxed time-dependent Schr\"odinger equation.}
\lectureframe{00:56:11}
\end{figure}

This is the generalized time-dependent Schr\"odinger equation. The familiar differential equation for a particle's wavefunction is one representation of this abstract state-vector equation. To use it, one must know \(H\) and the initial state.

For a finite interval, solve the differential equation. When \(H\) is time independent, repeatedly composing infinitesimal steps gives
\begin{equation}
U(t)=e^{-iHt}
\qquad(\hbar=1),
\end{equation}
which is the precise version of composing all the small steps. Time-dependent Hamiltonians require time ordering, a refinement we postpone.

\begin{classroomqa}
\audiencequestion{The expansion used an infinitesimal interval. What determines the state after a long time?}
\lecturerresponse{Apply the infinitesimal update repeatedly, or equivalently solve the differential equation. For time-independent \(H\), the product of the small steps becomes \(e^{-iHt/\hbar}\).}
\end{classroomqa}

\subsection{Restoring Planck's Constant}

The preceding equation used units where \(\hbar=1\). Dimensional analysis shows where \(\hbar\) belongs. The derivative has dimensions of inverse time, while \(H\) has dimensions of energy. Since \(\hbar\) has dimensions energy times time,
\begin{equation}
\boxed{i\hbar\frac{d}{dt}\ket{\psi(t)}=H\ket{\psi(t)}}.
\label{eq:l4-schrodinger}
\end{equation}
Equivalently, the infinitesimal operator is \(U(\epsilon)=I-i\epsilon H/\hbar+O(\epsilon^2)\), and the finite time-independent evolution is \(U(t)=e^{-iHt/\hbar}\).

\begin{classroomqa}
\audiencequestion{Where exactly does Planck's constant enter?}
\lecturerresponse{Planck's constant enters in either of the equivalent forms
\begin{align*}
i\hbar\,\partial_t\ket{\psi}&=H\ket{\psi},\\
\partial_t\ket{\psi}&=-\frac{i}{\hbar}H\ket{\psi}.
\end{align*}
Both have matching dimensions because \(\hbar\) has units of energy times time.}
\end{classroomqa}

The state evolves deterministically according to Equation~\eqref{eq:l4-schrodinger}; measurement outcomes remain statistical because they are sampled from the continuously changing state by Born's rule.

\begin{classroomqa}
\audiencequestion{Does the derivation depend on reversibility, and how can it survive a measurement that disturbs the state?}
\lecturerresponse{It depends on linear evolution and conservation of inner products. The equation governs an isolated system. A measuring apparatus interacts strongly enough to become part of the physical system and cannot be ignored. The combined system and apparatus must be treated quantum mechanically; their entanglement is precisely what must be understood before measurement can be analyzed.}
\end{classroomqa}

\section{Expectation Value Means Average}

To compare quantum dynamics with classical equations, we need the average value of an observable. The standard name \emph{expectation value} is potentially misleading. For a fair variable taking only \(+1\) and \(-1\), the average is zero although zero is never observed. A symmetric distribution can even vanish throughout a neighborhood of its mean. The mean, median, and mode coincide for some smooth, single-peaked distributions, but not in general.

If \(L\) can return values \(\lambda_i\) with probabilities \(P(\lambda_i)\), define
\begin{equation}
\overline{\lambda}=\sum_i\lambda_iP(\lambda_i).
\label{eq:l4-probability-average}
\end{equation}
This is the first moment of the distribution, not a promise about the result of one trial.

\begin{classroomqa}
\audiencequestion{Isn't the most frequent value what one should expect?}
\lecturerresponse{The most frequent value is the mode; the weighted average is the mean. They may coincide for a simple narrow distribution, but a two-peaked distribution can have its mean in a region of essentially zero probability. The word expectation refers to the mathematical mean, not necessarily to the most likely single outcome.}
\end{classroomqa}

Repeated trials connect the mathematical mean to data. Under the usual independence and finite-moment assumptions, laws of large numbers state probabilistic convergence of the sample average. They do not guarantee that every finite sequence lies close to the mean; an extraordinarily long run of identical spin outcomes remains possible, merely improbable. This is the same statistical subtlety that prevents the second law of thermodynamics from being an exceptionless microscopic prohibition.

\begin{classroomqa}
\audiencequestion{If a fair spin experiment happens to give a million ups in a row, is its average really up?}
\lecturerresponse{The sample average of that particular run is indeed up. The theoretical expectation remains zero. As the number of independent trials grows, the sample average approaches the theoretical mean with high probability, but probability theory never rules out every exceptional finite sequence.}
\end{classroomqa}

\subsection{The Operator Sandwich}

Expand a normalized state in an orthonormal eigenbasis of \(L\):
\begin{equation}
\ket{A}=\sum_i\alpha_i\ket{\lambda_i},
\qquad
\bra{A}=\sum_j\alpha_j^*\bra{\lambda_j}.
\end{equation}
Born's rule and Equation~\eqref{eq:l4-probability-average} give
\begin{equation}
\overline{\lambda}=\sum_i\lambda_i|\alpha_i|^2.
\end{equation}
Now calculate
\begin{align}
\bra{A}L\ket{A}
&=\sum_{i,j}\alpha_j^*\alpha_i
  \bra{\lambda_j}L\ket{\lambda_i}\\
&=\sum_{i,j}\alpha_j^*\alpha_i\lambda_i
  \braket{\lambda_j}{\lambda_i}\\
&=\sum_i\lambda_i|\alpha_i|^2.
\end{align}
The double sum collapses because
\(\braket{\lambda_j}{\lambda_i}=\delta_{ji}\). Hence
\begin{equation}
\boxed{\expect{L}_A=\bra{A}L\ket{A}}.
\label{eq:l4-expectation-sandwich}
\end{equation}

\begin{figure}[htbp]
\centering
\includegraphics[width=0.86\textwidth]{lecture_04_figure_06.png}
\caption{The eigenbasis expansion reduces the probabilistic average to the operator sandwich \(\bra{A}L\ket{A}\).}
\lectureframe{01:11:50}
\end{figure}

Angle brackets and an overbar will both be used for averages. When the state evolves,
\begin{equation}
\expect{L}_t=\bra{\psi(t)}L\ket{\psi(t)}.
\end{equation}
The subscript only reminds us that the state, and therefore the average, depends on time. In the present Schr\"odinger-picture discussion, \(L\) itself is fixed unless explicit time dependence is stated.

\begin{classroomqa}
\audiencequestion{Should the \(t\) be written outside the brackets, and does an overbar mean a sample mean?}
\lecturerresponse{Either notation can be annotated to show time dependence, but the standard quantum notation is \(\expect{L}_t\), or simply \(\expect{L}\) with the time understood from \(\ket{\psi(t)}\). Here an overbar denotes the theoretical probability average, the first moment \(\sum_i\lambda_iP(\lambda_i)\), not specifically a finite-data sample mean.}
\end{classroomqa}

\section{How Averages Evolve}

Differentiate the sandwich with the product rule:
\begin{equation}
\frac{d}{dt}\expect{L}
=
\left(\frac{d}{dt}\bra{\psi}\right)L\ket{\psi}
+
\bra{\psi}L\left(\frac{d}{dt}\ket{\psi}\right).
\label{eq:l4-product-derivative}
\end{equation}
In \(\hbar=1\) units, the ket equation and its Hermitian conjugate are
\begin{align}
\frac{d}{dt}\ket{\psi}&=-iH\ket{\psi},\\
\frac{d}{dt}\bra{\psi}&=+i\bra{\psi}H.
\end{align}
The sign of \(i\) changes under complex conjugation, and \(H^\dagger=H\). Substitution into Equation~\eqref{eq:l4-product-derivative} gives
\begin{align}
\frac{d}{dt}\expect{L}
&=i\bra{\psi}HL\ket{\psi}
-i\bra{\psi}LH\ket{\psi}\\
&=i\bra{\psi}(HL-LH)\ket{\psi}.
\end{align}

Operators need not commute. Their order matters just as matrix order matters. Define
\begin{equation}
[H,L]=HL-LH.
\end{equation}
Then
\begin{equation}
\boxed{\frac{d}{dt}\expect{L}=i\expect{[H,L]}}
\qquad(\hbar=1).
\label{eq:l4-average-evolution-natural}
\end{equation}

\begin{figure}[htbp]
\centering
\includegraphics[width=0.86\textwidth]{lecture_04_figure_07.png}
\caption{The time derivative of an average is governed by the expectation value of a commutator.}
\lectureframe{01:22:22}
\end{figure}

Restoring dimensions,
\begin{equation}
\boxed{\frac{d}{dt}\expect{L}
=\frac{i}{\hbar}\expect{[H,L]}
=-\frac{i}{\hbar}\expect{[L,H]}}.
\label{eq:l4-average-evolution}
\end{equation}
If \(L\) itself depends explicitly on time, the general equation includes the additional term \(\expect{\partial L/\partial t}\). Here \(L\) is fixed.

\begin{classroomqa}
\audiencequestion{Is the sign in the commutator equation really plus \(i\)?}
\lecturerresponse{Yes, with \([H,L]=HL-LH\), the result is \(+(i/\hbar)\expect{[H,L]}\). Reversing the commutator reverses its sign, giving the equivalent form \(-(i/\hbar)\expect{[L,H]}\).}
\end{classroomqa}

\begin{classroomqa}
\audiencequestion{Was the bra equation obtained by complex conjugation or Hermitian conjugation, and why does the derivative remain on the left?}
\lecturerresponse{It is Hermitian conjugation. The derivative of a ket is a ket, and the dual of that derivative is the derivative of the corresponding bra. One may see this directly from the difference quotient: differences of bras are bras, just as differences of kets are kets. Time is a parameter here, not an operator whose order must be reversed.}
\end{classroomqa}

\section{The Classical Shadow}

In classical Hamiltonian mechanics, a phase-space function \(L(q,p)\) with no explicit time dependence obeys
\begin{equation}
\frac{dL}{dt}=\{L,H\}_{\mathrm{PB}},
\label{eq:l4-poisson-evolution}
\end{equation}
where \(\{\cdot,\cdot\}_{\mathrm{PB}}\) is the Poisson bracket. Compare this with Equation~\eqref{eq:l4-average-evolution}. Both brackets are antisymmetric:
\begin{align}
\{L,H\}_{\mathrm{PB}}&=-\{H,L\}_{\mathrm{PB}},\\
[L,H]&=-[H,L].
\end{align}
Both also act as derivatives on a product:
\begin{align}
\{LM,H\}_{\mathrm{PB}}
&=\{L,H\}_{\mathrm{PB}}M+L\{M,H\}_{\mathrm{PB}},\\
[LM,H]
&=L[M,H]+[L,H]M.
\label{eq:l4-commutator-product}
\end{align}
The commutator identity follows by adding and subtracting \(LHM\):
\begin{align}
LMH-HLM
&=(LMH-LHM)\\
&\quad +(LHM-HLM)\\
&=L(MH-HM)\\
&\quad +(LH-HL)M.
\end{align}

The parallel suggests the correspondence
\begin{equation}
[A,B]
\quad\longleftrightarrow\quad
i\hbar\{A,B\}_{\mathrm{PB}},
\label{eq:l4-dirac-correspondence}
\end{equation}
or equivalently
\begin{equation}
\{A,B\}_{\mathrm{PB}}
\quad\longleftrightarrow\quad
-\frac{i}{\hbar}[A,B].
\end{equation}

\begin{figure}[htbp]
\centering
\includegraphics[width=0.86\textwidth]{lecture_04_figure_08.png}
\caption{The classical Poisson-bracket equation and the quantum commutator correspondence, with \(\hbar\) restored.}
\lectureframe{01:32:32}
\end{figure}

Dirac recognized this connection and used it as a route into quantum mechanics. It must nevertheless be handled as a correspondence, not a universal substitution rule: operator-ordering ambiguities can appear, and the classical limit must be derived from quantum dynamics. Conceptually, quantum mechanics comes first; classical mechanics is its approximation for suitable states and scales. Commutators precede Poisson brackets, even though the historical analogy helps us discover the relation.

\begin{classroomqa}
\audiencequestion{Why call \(H\) energy, and why infer energy conservation before imposing a symmetry principle?}
\lecturerresponse{At this stage \(H\) is rigorously the Hermitian generator of time translations. Calling it energy is a physical identification that must be checked where a classical limit is known. For large or heavy systems in suitable states, the quantum averages must follow the classical equations, and the quantum \(H\) must connect to the classical Hamiltonian. The conservation statement below assumes that this generator has no explicit time dependence.}
\end{classroomqa}

\section{Average Energy Is Conserved}

Take \(L=H\) in Equation~\eqref{eq:l4-average-evolution}, with a time-independent Hamiltonian:
\begin{equation}
\frac{d}{dt}\expect{H}
=\frac{i}{\hbar}\expect{[H,H]}.
\end{equation}
Every operator commutes with itself,
\begin{equation}
[H,H]=HH-HH=0,
\end{equation}
so
\begin{equation}
\boxed{\frac{d}{dt}\expect{H}=0}.
\label{eq:l4-energy-conservation}
\end{equation}

\begin{figure}[htbp]
\centering
\includegraphics[width=0.86\textwidth]{lecture_04_figure_09.png}
\caption{Average-energy conservation follows immediately because the Hamiltonian commutes with itself.}
\lectureframe{01:39:30}
\end{figure}

This proves a deliberately modest statement: the average energy is constant. A stronger statement about the entire energy distribution will follow from the same dynamics. If \(H\) depends explicitly on time, the extra explicit derivative need not vanish and the driven system can exchange energy with whatever changes its parameters.

Two threads remain open. The arbitrary-axis spin probability calculation is still waiting, and the state-update problem has not yet been derived. Between interventions an isolated state evolves unitarily. During a measurement, system and apparatus interact and become entangled.

\begin{classroomqa}
\audiencequestion{What happened to the promised probability for a spin measured along an arbitrary axis?}
\lecturerresponse{The general dynamics took the available time. The spin calculation remains the next concrete application and will be resumed after this structural framework.}
\end{classroomqa}

\begin{classroomqa}
\audiencequestion{How can the postulates be complete without a collapse postulate?}
\lecturerresponse{For now collapse is provisional language. A measurement establishes entanglement between the system and an apparatus. Until composite systems and entanglement have been developed, the measurement process cannot be modeled honestly. That is the next conceptual step.}
\end{classroomqa}
